\textbf{Phần 1. Trắc nghiệm lựa chọn}
\setcounterpageref{ex}{0}

\begin{ex} %[5D7N43]
Bạn Nam mua $3$ quyển vở giá $8\,500$ đồng/quyển và $2$ cây bút giá $5\,000$ đồng/cây. Biểu thức tính tổng số tiền bằng máy tính cầm tay là
\choice
{$3 \times 8\,500 \times 2 \times 5\,000$}
{$3 + 8\,500 + 2 + 5\,000$}
{\True $3 \times 8\,500 + 2 \times 5\,000$}
{$(3 + 2) \times (8\,500 + 5\,000)$}
\loigiai{
Tổng số tiền mua vở và bút là: $3 \times 8\,500 + 2 \times 5\,000$.\\
\textbf{Đáp án của câu hỏi là C}
}
\end{ex}

\begin{ex} %[5D7N43]
Một giỏ hoa quả gồm $1{,}5\text{ kg}$ táo với giá $40\,000$ đồng/kg. Số tiền mua táo tính được là
\choice
{$50\,000$ đồng}
{$70\,000$ đồng}
{$80\,000$ đồng}
{\True $60\,000$ đồng}
\loigiai{
Số tiền: $1{,}5 \times 40\,000 = 60\,000$ (đồng).\\
\textbf{Đáp án của câu hỏi là D}
}
\end{ex}

\begin{ex} %[5D7H43]
Mẹ đưa cho Mai $200\,000$ đồng đi siêu thị. Mai mua hết $146\,500$ đồng. Sử dụng máy tính cầm tay tính số tiền thừa Mai nhận lại là
\choice
{$63\,500$ đồng}
{$54\,500$ đồng}
{\True $53\,500$ đồng}
{$64\,500$ đồng}
\loigiai{
Số tiền thừa: $200\,000 - 146\,500 = 53\,500$ (đồng).\\
\textbf{Đáp án của câu hỏi là C}
}
\end{ex}

\begin{ex} %[5D7H43]
Một cửa hàng giảm giá $10\%$ cho hóa đơn mua hàng trên $500\,000$ đồng. Bác An mua hàng hết $650\,000$ đồng. Số tiền bác An được giảm là
\choice
{$50\,000$ đồng}
{\True $65\,000$ đồng}
{$60\,000$ đồng}
{$70\,000$ đồng}
\loigiai{
Số tiền được giảm là: $650\,000 \times 10\% = 65\,000$ (đồng).\\
\textbf{Đáp án của câu hỏi là B}
}
\end{ex}

\begin{ex} %[5D7H43]
Một mảnh vườn hình chữ nhật có chiều dài $24{,}5\text{ m}$ và chiều rộng $16\text{ m}$. Sử dụng máy tính cầm tay tính diện tích mảnh vườn là
\choice
{$384\text{ m}^2$}
{$400\text{ m}^2$}
{$412\text{ m}^2$}
{\True $392\text{ m}^2$}
\loigiai{
Diện tích: $24{,}5 \times 16 = 392\text{ (m}^2\text{)}$.\\
\textbf{Đáp án của câu hỏi là D}
}
\end{ex}

\begin{ex} %[5D7H43]
Bác Ba thu hoạch được $1{,}2$ tấn thóc. Bác để lại $15\%$ để ăn, số còn lại đem bán. Dùng máy tính cầm tay tính khối lượng thóc bác Ba đem bán là
\choice
{$180\text{ kg}$}
{$1\,000\text{ kg}$}
{\True $1\,020\text{ kg}$}
{$980\text{ kg}$}
\loigiai{
Đổi $1{,}2$ tấn $= 1\,200\text{ kg}$.\\
Khối lượng thóc đem bán: $1\,200 \times (100\% - 15\%) = 1\,020\text{ (kg)}$.\\
\textbf{Đáp án của câu hỏi là C}
}
\end{ex}

\begin{ex} %[5D7V43]
Lan mua $5$ quyển sách tham khảo giá $35\,000$ đồng/quyển và $1$ chiếc cặp sách giá $185\,000$ đồng. Nhân dịp năm học mới cửa hàng giảm $10\%$ trên tổng hóa đơn. Lan phải trả số tiền là
\choice
{$360\,000$ đồng}
{$330\,000$ đồng}
{$340\,000$ đồng}
{\True $324\,000$ đồng}
\loigiai{
Tổng hóa đơn trước giảm: $5 \times 35\,000 + 185\,000 = 175\,000 + 185\,000 = 360\,000$ (đồng).\\
Số tiền phải trả sau khi giảm $10\%$: $360\,000 \times 90\% = 324\,000$ (đồng).\\
\textbf{Đáp án của câu hỏi là D}
}
\end{ex}

\begin{ex} %[5D7V43]
Một gia đình có thu nhập hàng tháng là $25\,000\,000$ đồng. Gia đình dành $45\%$ chi phí cho ăn uống, $25\%$ cho học tập, $15\%$ cho các chi tiêu khác, phần còn lại để tiết kiệm. Số tiền tiết kiệm mỗi tháng là
\choice
{\True $3\,750\,000$ đồng}
{$3\,500\,000$ đồng}
{$4\,000\,000$ đồng}
{$4\,250\,000$ đồng}
\loigiai{
Tỉ lệ tiết kiệm: $100\% - (45\% + 25\% + 15\%) = 15\%$.\\
Số tiền tiết kiệm: $25\,000\,000 \times 15\% = 3\,750\,000$ (đồng).\\
\textbf{Đáp án của câu hỏi là A}
}
\end{ex}

\begin{ex} %[5D7V43]
Một cửa hàng điện máy bán một máy giặt với giá $7\,800\,000$ đồng. Khách hàng thanh toán qua thẻ ngân hàng được hoàn tiền $5\%$. Số tiền thực tế khách hàng phải trả là
\choice
{$7\,500\,000$ đồng}
{\True $7\,410\,000$ đồng}
{$7\,400\,000$ đồng}
{$7\,390\,000$ đồng}
\loigiai{
Số tiền được hoàn: $7\,800\,000 \times 5\% = 390\,000$ (đồng).\\
Số tiền thực tế phải trả: $7\,800\,000 - 390\,000 = 7\,410\,000$ (đồng).\\
\textbf{Đáp án của câu hỏi là B}
}
\end{ex}

\begin{ex} %[5D7V43]
Bác nông dân nuôi $250$ con gà và $150$ con vịt. Sau một đợt dịch bệnh, tỉ lệ sống sót của gà là $96\%$, của vịt là $94\%$. Tổng số gà và vịt còn sống sót là
\choice
{\True $381$ con}
{$375$ con}
{$385$ con}
{$390$ con}
\loigiai{
Số gà sống sót: $250 \times 96\% = 240$ (con).\\
Số vịt sống sót: $150 \times 94\% = 141$ (con).\\
Tổng số con sống sót: $240 + 141 = 381$ (con).\\
\textbf{Đáp án của câu hỏi là A}
}
\end{ex}

\begin{ex} %[5D7V43]
Một hình hộp chữ nhật có chiều dài $1{,}2\text{ m}$, chiều rộng $0{,}8\text{ m}$ và chiều cao $0{,}5\text{ m}$. Sử dụng máy tính cầm tay tính thể tích của hình hộp này theo đơn vị lít ($1\text{ dm}^3 = 1\text{ l}$).
\choice
{$48\text{ l}$}
{\True $480\text{ l}$}
{$4\,800\text{ l}$}
{$4800\text{ l}$}
\loigiai{
Thể tích: $1{,}2 \times 0{,}8 \times 0{,}5 = 0{,}48\text{ (m}^3\text{)} = 480\text{ dm}^3 = 480\text{ l}$.\\
\textbf{Đáp án của câu hỏi là B}
}
\end{ex}

\begin{ex} %[5D7C43]
Một siêu thị có chương trình khuyến mãi: Mua $3$ hộp sữa tặng $1$ hộp sữa cùng loại. Mỗi hộp sữa có giá $28\,000$ đồng. Mẹ An cần $14$ hộp sữa thì phải trả ít nhất bao nhiêu tiền?
\choice
{\True $308\,000$ đồng}
{$336\,000$ đồng}
{$292\,000$ đồng}
{$392\,000$ đồng}
\loigiai{
Mua 3 tặng 1 nghĩa là mua 4 hộp chỉ trả tiền 3 hộp ($3 \times 28\,000 = 84\,000$ đồng).\\
Ta có: $14 = 4 \times 3 + 2$.\\
Số hộp phải trả tiền: $3 \times 3 + 2 = 11$ hộp.\\
Tổng số tiền phải trả: $11 \times 28\,000 = 308\,000$ (đồng).\\
\textbf{Đáp án của câu hỏi là A}
}
\end{ex}

\textbf{Phần 2. Trắc nghiệm đúng sai}
\setcounterpageref{ex}{0}

\begin{ex} %[5D7H43]
Bạn Hoa đi mua đồ dùng học tập gồm: $5$ cuốn tập giá $12\,000$ đồng/cuốn, $3$ cây bút bi giá $6\,000$ đồng/cây và $1$ hộp bút giá $45\,000$ đồng. Xét tính đúng sai:
\choiceTF
{\True Tổng số tiền mua $5$ cuốn tập là $60\,000$ đồng.}
{\True Tổng số tiền mua $3$ cây bút bi là $18\,000$ đồng.}
{Tổng tiền hóa đơn khi chưa giảm giá là $133\,000$ đồng.}
{\True Nếu được giảm giá $10\%$, Hoa phải trả $110\,700$ đồng.}
\loigiai{
\begin{itemize}
    \item Tiền tập: $5 \times 12\,000 = 60\,000$ đồng. Đúng.
    \item Tiền bút bi: $3 \times 6\,000 = 18\,000$ đồng. Đúng.
    \item Tổng hóa đơn: $60\,000 + 18\,000 + 45\,000 = 123\,000$ đồng $\neq 133\,000$ đồng. Sai.
    \item Tiền sau giảm 10\%: $123\,000 \times 0{,}9 = 110\,700$ đồng. Đúng.
\end{itemize}
\begin{itemize}
    \item a) \textbf{Đúng}.
    \item b) \textbf{Đúng}.
    \item c) \textbf{Sai}.
    \item d) \textbf{Đúng}.
\end{itemize}
}
\end{ex}

\begin{ex} %[5D7V43]
Một xí nghiệp may có $120$ công nhân. Mỗi công nhân trung bình may được $15$ chiếc áo một ngày. Mỗi chiếc áo bán với giá $150\,000$ đồng. Xét tính đúng sai:
\choiceTF
{\True Trong một ngày, xí nghiệp may được $1\,800$ chiếc áo.}
{\True Doanh thu từ số áo may trong một ngày là $270$ triệu đồng.}
{Nếu tăng năng suất thêm $10\%$ thì mỗi ngày may được $1\,900$ chiếc áo.}
{\True Sau khi tăng năng suất $10\%$, mỗi ngày xí nghiệp may được $1\,980$ chiếc áo.}
\loigiai{
\begin{itemize}
    \item Số áo 1 ngày: $120 \times 15 = 1\,800$ chiếc. Đúng.
    \item Doanh thu: $1\,800 \times 150\,000 = 270\,000\,000$ đồng $= 270$ triệu đồng. Đúng.
    \item Tăng 10\%: $1\,800 \times 1{,}1 = 1\,980$ chiếc $\neq 1\,900$ chiếc. Ý 3 Sai, Ý 4 Đúng.
\end{itemize}
\begin{itemize}
    \item a) \textbf{Đúng}.
    \item b) \textbf{Đúng}.
    \item c) \textbf{Sai}.
    \item d) \textbf{Đúng}.
\end{itemize}
}
\end{ex}

\textbf{Phần 3. Trả lời ngắn}
\setcounterpageref{ex}{0}

\begin{ex} %[5D7H43]
Sử dụng máy tính cầm tay tính tổng số tiền mua hàng: $4$ hộp bánh giá $45\,000$ đồng/hộp và $6$ chai nước ngọt giá $12\,500$ đồng/chai. (Ghi kết quả dưới dạng số nghìn đồng, ví dụ: 255).
\shortans{255}
\loigiai{
\textbf{Đáp số:} $255$\\
$4 \times 45\,000 + 6 \times 12\,500 = 180\,000 + 75\,000 = 255\,000$ đồng $= 255$ nghìn đồng.
}
\end{ex}

\begin{ex} %[5D7H43]
Một người đi chợ mua $2{,}5\text{ kg}$ thịt lợn với giá $130\,000$ đồng/kg và $3\text{ kg}$ cá với giá $85\,000$ đồng/kg. Tính tổng số tiền người đó phải trả theo đơn vị nghìn đồng.
\shortans{580}
\loigiai{
\textbf{Đáp số:} $580$\\
$2{,}5 \times 130 + 3 \times 85 = 325 + 255 = 580$ nghìn đồng.
}
\end{ex}

\begin{ex} %[5D7V43]
Một mảnh đất hình chữ nhật có chiều dài $32{,}5\text{ m}$, chiều rộng $18\text{ m}$. Người ta dùng $20\%$ diện tích để xây nhà. Tính diện tích xây nhà theo mét vuông.
\shortans{117}
\loigiai{
\textbf{Đáp số:} $117$\\
Diện tích mảnh đất: $32{,}5 \times 18 = 585\text{ m}^2$.\\
Diện tích xây nhà: $585 \times 20\% = 117\text{ m}^2$.
}
\end{ex}

\begin{ex} %[5D7V43]
Một gia đình sử dụng hết $185\text{ kWh}$ điện trong tháng. Giá $50\text{ kWh}$ đầu là $1\,800$ đồng/kWh, $50\text{ kWh}$ tiếp theo là $1\,900$ đồng/kWh, số còn lại là $2\,100$ đồng/kWh. Tính tổng số tiền điện phải trả (chưa tính thuế) theo đơn vị nghìn đồng.
\shortans{363,5}
\loigiai{
\textbf{Đáp số:} $363{,}5$\\
Tiền điện: $50 \times 1{,}8 + 50 \times 1{,}9 + 85 \times 2{,}1 = 90 + 95 + 178{,}5 = 363{,}5$ nghìn đồng.
}
\end{ex}

\begin{ex} %[5D7V43]
Một cửa hàng quần áo niêm yết giá một chiếc áo khoác là $650\,000$ đồng. Trong chương trình tri ân, áo được giảm giá $20\%$, khách quen lại được giảm thêm $5\%$ trên giá đã giảm. Hỏi khách quen mua áo khoác đó với giá bao nhiêu nghìn đồng?
\shortans{494}
\loigiai{
\textbf{Đáp số:} $494$\\
Giá sau giảm 20\%: $650 \times 0{,}8 = 520$ nghìn đồng.\\
Giá khách quen: $520 \times 0{,}95 = 494$ nghìn đồng.
}
\end{ex}

\begin{ex} %[5D7C43]
Một bác nông dân bán $450\text{ kg}$ bưởi với giá $35\,000$ đồng/kg. Bác dành $40\%$ số tiền thu được để mua phân bón, $\dfrac{1}{3}$ số tiền còn lại để mua giống cây mới, số tiền thừa bác gửi vào quỹ tiết kiệm. Hỏi bác gửi vào quỹ tiết kiệm bao nhiêu nghìn đồng?
\shortans{6300}
\loigiai{
\textbf{Đáp số:} $6300$\\
Tổng tiền: $450 \times 35 = 15\,750$ nghìn đồng.\\
Mua phân bón (40\%): $15\,750 \times 0{,}4 = 6\,300$ nghìn đồng.\\
Còn lại: $15\,750 - 6\,300 = 9\,450$ nghìn đồng.\\
Mua giống (1/3 còn lại): $9\,450 : 3 = 3\,150$ nghìn đồng.\\
Gửi tiết kiệm: $9\,450 - 3\,150 = 6\,300$ nghìn đồng.
}
\end{ex}
